\documentclass{article}
\usepackage[french]{babel}
\usepackage[T1]{fontenc}
\usepackage[utf8]{inputenc}
\usepackage{graphicx}
\usepackage{mathtools}
\usepackage{listingsutf8}
\usepackage{amsfonts}
\usepackage{amssymb}
\lstset{%
    language=Python,
    numbers=left,
    xleftmargin=2em,
    inputencoding=utf8,
    basicstyle=\ttfamily\small,
    keywordstyle=\bfseries,
    showstringspaces=false,
    literate=
        {é}{{\'e}}{1}%
        {è}{{\`e}}{1}%
        {à}{{\`a}}{1}%
        {â}{{\^a}}{1}%
        {ç}{{\c{c}}}{1}%
        {œ}{{\oe}}{1}%
        {ù}{{\`u}}{1}%
        {É}{{\'E}}{1}%
        {È}{{\`E}}{1}%
        {À}{{\`A}}{1}%
        {Ç}{{\c{C}}}{1}%
        {Œ}{{\OE}}{1}%
        {Ê}{{\^E}}{1}%
        {ê}{{\^e}}{1}%
        {î}{{\^i}}{1}%
        {ï}{{\"i}}{1}%
        {ô}{{\^o}}{1}%
        {û}{{\^u}}{1}%
}
\usepackage{fullpage}
\usepackage[margin=1.5cm]{geometry}
\usepackage{setspace}
\usepackage{todonotes}
\usepackage{amsthm}

\newtheoremstyle{exostyle}
{10pt}
{10pt}
{}
{}
{\scshape\bfseries\large}
{\hfill\vspace{5pt}\newline}
{0pt}
{\hfill\thmname{#1}\thmnumber{ #2} -- \thmnote{ #3}}

\newtheoremstyle{partiestyle}
{1em}
{1em}
{}
{}
{\bfseries}
{\vspace{.5em}\newline}
{0em}
{\thmnumber{#2} \thmnote{ #3}}

\newtheoremstyle{questionstyle}
{.5em}
{.5em}
{}
{}
{\bfseries}
{}
{0em}
{Question \thmnumber{#2 }}

\theoremstyle{exostyle}
\newtheorem{exo}{Exercice}

\theoremstyle{partiestyle}
\newtheorem{partie}{}[exo]

\theoremstyle{questionstyle}
\newtheorem{question}{Question}[exo]
\newtheorem{questionpartie}{Question}[partie]

\title{Examen terminal: Programmation}
\author{L1 MPCI}
\date{28 mai 2026 - Durée: 2h}

\begin{document}

\maketitle

\begin{center}
{\em On rappelle qu'aucun document, ni équipement électrique ou électronique n'est autorisé.

{\bf Bien que} une Game Boy Color pourra être utilisée sans le son si vous avez besoin de patienter.}
\end{center}

\vspace*{1cm}

\paragraph{Objectif de l'examen} Cet examen évalue votre capacité à concevoir et implémenter un programme orienté objet en Python, et à appliquer le \textit{design pattern} MVC (Modèle-Vue-Contrôleur). À travers une application de gestion de bibliothèque, vous serez amenés à définir des classes, à comprendre la séparation des responsabilités, puis à structurer un programme complet en trois composants distincts.

\paragraph{Les exercices de cet examen} :
\begin{itemize}
\item sont au nombre de 3 ;
\item les exercices ne sont pas indépendants mais peuvent être menés en parallèle ;
\item sont chacun constitués de plusieurs parties de difficultés croissantes ;
\end{itemize}

\paragraph{}{\large \sc Rendez \underline{des copies séparées pour chaque exercice}, ceci vous permettra de d'y revenir au cours de l'examen sans perdre le correcteur.}

\vfill
\begin{center}
    {\small \em Sujet conçu par Claude (et) François.}
\end{center}

\clearpage

% ============================================================
\begin{exo}[Programmation orientée objet]

Le contexte de cet exercice est une application de gestion d'une petite bibliothèque personnelle.

    \begin{partie}[La classe \texttt{Livre}]

On vous donne la classe suivante~:

\begin{lstlisting}
class Livre:
    def __init__(self, titre, auteur):
        self.titre = titre
        self.auteur = auteur
        self._disponible = True
\end{lstlisting}

        \begin{questionpartie}
            Qu'appelle-t-on un \textit{attribut}~? Listez les attributs d'instance de la classe \texttt{Livre} et donnez leur type.
        \end{questionpartie}

        \begin{questionpartie}
            Pourquoi l'attribut \texttt{\_disponible} est-il préfixé d'un tiret bas~? Ajoutez une méthode de nom \texttt{est\_disponible} qui rend sa valeur.
        \end{questionpartie}

        \begin{questionpartie}
            Complétez la méthode suivante. Elle doit retourner \texttt{True} si le livre était disponible (et le marquer comme emprunté), \texttt{False} sinon~:

\begin{lstlisting}
def emprunter(self):
    ...
\end{lstlisting}
        \end{questionpartie}

        \begin{questionpartie}
            Écrivez la méthode \texttt{\_\_str\_\_}. Elle doit retourner une chaîne du type~:
            \begin{center}
                \texttt{"Trilogie spinoziste (Jean-Bernard Pouy) [disponible]"}
            \end{center}
            en indiquant \texttt{[disponible]} ou \texttt{[emprunté]} selon l'état du livre.
        \end{questionpartie}

    \end{partie}

    \begin{partie}[La classe \texttt{Bibliothèque}]

On vous donne le début de la classe suivante~:

\begin{lstlisting}
class Bibliothèque:
    def __init__(self):
        self._livres = []
\end{lstlisting}

        \begin{questionpartie}
            Ajoutez une méthode \texttt{Bibliothèque.ajouter(livre)} qui ajoute un livre à une bibliothèque et une méthode \texttt{Bibliothèque.livres\_disponibles()} qui retourne la liste des livres actuellement disponibles.
        \end{questionpartie}

        \begin{questionpartie}
            Ajoutez une méthode \texttt{Bibliothèque.chercher(titre)} qui retourne le \texttt{Livre} dont le titre correspond, ou \texttt{None} si introuvable.

            Puis, en utilisant \texttt{chercher}, ajoutez une méthode \texttt{Bibliothèque.emprunter(titre)} qui tente d'emprunter le livre portant ce titre et retourne \texttt{True} en cas de succès.
        \end{questionpartie}

    \end{partie}

    \begin{partie}[Héritage : le livre numérique]

        On souhaite modéliser les livres numériques. Contrairement à un livre physique, un livre numérique peut toujours être emprunté : il existe en nombre illimité de copies.

        \begin{questionpartie}
            Qu'appelle-t-on l'héritage en programmation orientée objet~? Quelle relation établit-il entre une classe parente et une classe fille~?
        \end{questionpartie}

        \begin{questionpartie}
            Écrivez une classe \texttt{LivreNumérique} qui hérite de \texttt{Livre}. Son constructeur prendra les mêmes paramètres que celui de \texttt{Livre} et appellera le constructeur parent à l'aide de \texttt{super()}.
        \end{questionpartie}

        \begin{questionpartie}
            Surchargez la méthode \texttt{emprunter} dans \texttt{LivreNumérique} de façon à ce qu'elle retourne toujours \texttt{True}, sans modifier l'attribut \texttt{\_disponible}.
        \end{questionpartie}

        \begin{questionpartie}
            Surchargez \texttt{\_\_str\_\_} pour retourner une chaîne indiquant que le livre est numérique, par exemple~:
            \begin{center}
                \texttt{"Trilogie spinoziste (Jean-Bernard Pouy) [numérique]"}
            \end{center}
        \end{questionpartie}

        \begin{questionpartie}
            Sans modifier la classe \texttt{Bibliothèque}, peut-on lui ajouter un \texttt{LivreNumérique} via \texttt{ajouter}~? La méthode \texttt{emprunter} de \texttt{Bibliothèque} fonctionnera-t-elle correctement avec ce livre~? Quel principe de la programmation orientée objet illustre ce comportement~?
        \end{questionpartie}

    \end{partie}

    \begin{partie}[Modélisation UML]

        \begin{questionpartie}
            Dessinez la boîte UML de la classe \texttt{Livre}.
        \end{questionpartie}

        \begin{questionpartie}
            Ajoutez \texttt{LivreNumérique} au diagramme.
        \end{questionpartie}

        \begin{questionpartie}
            Ajoutez \texttt{Bibliothèque} au diagramme.
        \end{questionpartie}

    \end{partie}
\end{exo}

% ============================================================
\begin{exo}[Séparation des responsabilités]

On vous donne le code suivant d'une application de bibliothèque écrite sans aucune architecture particulière~:

\begin{lstlisting}
def gerer_bibliotheque():
    livres = []  # liste de dict {"titre": str, "auteur": str, "dispo": bool}

    while True:
        print("1. Livres disponibles  2. Emprunter  0. Quitter")
        choix = input("Votre choix : ")                      # (a)

        if choix == "1":
            dispo = [l for l in livres if l["dispo"]]        # (b)
            for l in dispo:
                print(f"- {l['titre']} ({l['auteur']})")     # (c)
        elif choix == "2":
            titre = input("Titre : ")                        # (d)
            livre = next((l for l in livres                  # (e)
                          if l["titre"] == titre), None)
            if livre and livre["dispo"]:
                livre["dispo"] = False                       # (f)
                print("Emprunté.")
            else:
                print("Indisponible ou introuvable.")
        elif choix == "0":
            break
\end{lstlisting}

    \begin{partie}[Analyse]

        Un \textit{design pattern} (ou patron de conception) est une solution générale et réutilisable à un problème récurrent de conception logicielle. Il ne s'agit pas de code prêt à l'emploi, mais d'un modèle qui guide la structure d'un programme.

        \begin{questionpartie}
            Le \textit{design pattern} MVC découpe une application en trois composants~:
            \begin{itemize}
                \item le \textbf{Modèle} gère les données et les règles métier~;
                \item la \textbf{Vue} gère l'affichage et la saisie~;
                \item le \textbf{Contrôleur} coordonne le modèle et la vue.
            \end{itemize}
            Pour chacun des éléments \texttt{(a)} à \texttt{(f)} du code ci-dessus, indiquez s'il relève du Modèle, de la Vue ou du Contrôleur. Justifiez brièvement.
        \end{questionpartie}

        \begin{questionpartie}
            Donnez deux inconvénients concrets de ce code monolithique (c'est-à-dire où toute la logique, l'affichage et la gestion des données sont mélangés dans une seule et même fonction) par rapport à une architecture MVC.
        \end{questionpartie}

        \begin{questionpartie}
            Un développeur souhaite remplacer l'interface en ligne de commande par une interface web. Quelles parties du code serait-il obligé de réécrire~? Justifiez en vous appuyant sur le MVC.
        \end{questionpartie}
    \end{partie}

    \begin{partie}[Refactoring par composition]

        \begin{questionpartie}
            Quelle est la différence entre héritage et composition~? En vous appuyant sur les classes \texttt{Livre} et \texttt{Bibliothèque}, donnez un exemple de chacune des deux relations.
        \end{questionpartie}

        \begin{questionpartie}
            On souhaite transformer la fonction \texttt{gerer\_bibliotheque} en une classe \texttt{Application} qui utilise par composition une instance de \texttt{Bibliothèque}. Écrivez le constructeur de \texttt{Application} qui crée et stocke une \texttt{Bibliothèque} vide comme attribut.
        \end{questionpartie}

        \begin{questionpartie}
            Transformez le bloc \texttt{elif choix == "2":} (lignes 13 à 20) en une méthode \texttt{emprunter(self)} de la classe \texttt{Application}, en utilisant l'attribut \texttt{Bibliothèque} créé à la question précédente.
        \end{questionpartie}

    \end{partie}

    \begin{partie}[Modélisation UML]

        \begin{questionpartie}
            Dessinez le diagramme de classes UML de la classe \texttt{Application} (définie en partie~2) ainsi que son lien avec la classe \texttt{Bibliothèque}.
        \end{questionpartie}

        \begin{questionpartie}
            Dans l'architecture MVC, le Contrôleur dépend du Modèle et de la Vue. Représentez schématiquement en UML la différence entre~:
            \begin{itemize}
                \item un Contrôleur qui \textit{crée} lui-même le Modèle et la Vue dans son constructeur~;
                \item un Contrôleur qui \textit{reçoit} le Modèle et la Vue en paramètre de son constructeur.
            \end{itemize}
            Quelle relation UML correspond à chaque cas~? Quel avantage apporte le second cas~?
        \end{questionpartie}

    \end{partie}
\end{exo}

% ============================================================
\begin{exo}[Implémentation du \textit{design pattern} MVC]

On reprend les classes \texttt{Livre} et \texttt{Bibliothèque} de l'exercice~1.

    \begin{partie}[La Vue]

        On vous donne le diagramme UML suivant~:

        \begin{center}
        \renewcommand{\arraystretch}{1.0}
        \begin{tabular}{|p{7cm}|}
        \hline
        \multicolumn{1}{|c|}{\textbf{VueTerminal}} \\
        \hline
        \multicolumn{1}{|c|}{\phantom{(aucun attribut)}} \\
        \hline
        \texttt{+ afficher\_livres(livres\,:\,list)} \\
        \texttt{+ afficher\_message(message\,:\,str)} \\
        \texttt{+ demander\_titre() : str} \\
        \hline
        \end{tabular}
        \end{center}

        \begin{questionpartie}
            Implémentez en Python la classe \texttt{VueTerminal}. \texttt{afficher\_livres} affichera chaque livre sur une ligne en utilisant \texttt{str(livre)}~; si la liste est vide, elle affichera \texttt{"Aucun livre disponible."}.
        \end{questionpartie}

    \end{partie}

    \begin{partie}[Le Contrôleur]

        On vous donne le diagramme UML suivant~:

        \begin{center}
        \renewcommand{\arraystretch}{1.0}
        \begin{tabular}{|p{9cm}|}
        \hline
        \multicolumn{1}{|c|}{\textbf{Contrôleur}} \\
        \hline
        \texttt{- \_modele\,:\,Bibliothèque} \\
        \texttt{- \_vue\,:\,VueTerminal} \\
        \hline
        \texttt{+ \_\_init\_\_(modele\,:\,Bibliothèque,\,vue\,:\,VueTerminal)} \\
        \texttt{+ afficher\_disponibles()} \\
        \texttt{+ emprunter()} \\
        \texttt{+ lancer()} \\
        \hline
        \end{tabular}
        \end{center}

        \begin{questionpartie}
            Implémentez \texttt{\_\_init\_\_}, \texttt{afficher\_disponibles} et \texttt{emprunter}.

            \texttt{emprunter} doit~: demander un titre via la vue, tenter l'emprunt via le modèle, puis afficher un message de succès ou d'échec via la vue.
        \end{questionpartie}

        \begin{questionpartie}
            Implémentez \texttt{lancer} qui affiche en boucle le menu suivant et appelle la méthode du contrôleur correspondante~:
            \begin{center}
                \texttt{1. Livres disponibles \quad 2. Emprunter \quad 0. Quitter}
            \end{center}
        \end{questionpartie}

    \end{partie}

    \begin{partie}[Extensibilité]

        \begin{questionpartie}
            Écrivez le code de lancement de l'application~: créez un modèle, ajoutez-y deux livres, créez une vue et un contrôleur, puis lancez l'application.
        \end{questionpartie}

        \begin{questionpartie}
            On souhaite remplacer \texttt{VueTerminal} par une classe \texttt{VueGraphique} sans modifier \texttt{Contrôleur}. Quelles méthodes \texttt{VueGraphique} doit-elle obligatoirement posséder~? Justifiez.
        \end{questionpartie}

    \end{partie}

    \begin{partie}[Modélisation UML]

        \begin{questionpartie}
            Dessinez le diagramme de classes UML complet de l'architecture MVC implémentée dans cet exercice. Faites apparaître \texttt{Contrôleur}, \texttt{Bibliothèque} et \texttt{VueTerminal} avec leurs attributs, leurs méthodes principales et les relations qui les unissent (nature de chaque relation et multiplicités).
        \end{questionpartie}

        \begin{questionpartie}
            Ajoutez \texttt{VueGraphique} au diagramme. Quelle relation entretient-elle avec \texttt{VueTerminal}~? Comment le diagramme montre-t-il que \texttt{Contrôleur} peut fonctionner indifféremment avec l'une ou l'autre~?
        \end{questionpartie}

    \end{partie}
\end{exo}

\end{document}
